% GENERATED FILE. Edit canonical lessons, then run npm run generate:books.
% canonical-source-hash: fnv1a64:f551211c36587ac1
% canonical-lessons: AR-C463-munakh, AR-C463-ihtibas-harari, AR-C463-ghilaf-jawwi, AR-C463-inbiath, AR-C463-karbun

\chapter{The climate}
\label{ch:ar-the-climate}
\hlchaptermodality{463}

\begin{chapteropening}
\textbf{By the end of this chapter:} \emph{I can say climate, global warming, the atmosphere, an emission and carbon.}

Complete the last lesson of chapter 463: I can say climate, global warming, the atmosphere, an emission and carbon.
\end{chapteropening}

\section[\texorpdfstring{\ar{مناخ} (munākh)}{مناخ (munākh)}]{\ar{مناخ} (munākh) — climate}

\label{lesson:AR-C463-munakh}



\begin{quote}
Before the new one: say the Arabic for to have room for, then the Arabic for to raise.
\end{quote}

\subsection*{You'll want to know: \ar{مناخ}}
\textbf{\ar{مناخ}} — \emph{munākh} — \textquotedblleft{}climate\textquotedblright{}.

\begin{quote}
\ar{مناخ} \ar{الصحراء} \ar{جاف} \ar{جدا}.
\end{quote}

\begin{quote}
\emph{munākh aṣ-ṣaḥrāʾ jāff jiddan.}
\end{quote}

\begin{quote}
The climate of the desert is very dry.
\end{quote}

\begin{cousinweb}
\textbf{\ar{مناخ}} (\emph{munākh}) is built on the native Arabic root \emph{n-w-kh}, \textquotedblleft{}to make a camel kneel\textquotedblright{} (Lane). A \emph{munākh} was a halting place where camels knelt; modern Arabic gave the old word the sense of climate.
\end{cousinweb}

\begin{grammarlens}[title={What you've built}]
One more word for getting around.
\end{grammarlens}

\subsection*{Your turn}
\begin{itemize}
\raggedright
  \item \emph{Say it:} \emph{munākh}
  \item \emph{Say it:} \emph{munākh}, once more
  \item \emph{Say it:} \emph{munākh aṣ-ṣaḥrāʾ jāff jiddan.} — The climate of the desert is very dry.
  \item \emph{Recall:} say the Arabic for to have room for, then the Arabic for to raise, then say \emph{munākh} again
\end{itemize}

\begin{tcolorbox}[breakable,colback=teal!4,colframe=teal!35!black,title={Before you move on}]
What does \ar{مناخ} mean? (\textquotedblleft{}Climate\textquotedblright{}.) Say it once more.
\end{tcolorbox}
\section[\texorpdfstring{\ar{احتباس} \ar{حراري} (iḥtibās ḥarārī)}{احتباس حراري (iḥtibās ḥarārī)}]{\ar{احتباس} \ar{حراري} (iḥtibās ḥarārī) — global warming}

\label{lesson:AR-C463-ihtibas-harari}



\begin{quote}
Before the new one: say the Arabic for to raise, then the Arabic for climate.
\end{quote}

\subsection*{You'll want to know: \ar{احتباس} \ar{حراري}}
\textbf{\ar{احتباس} \ar{حراري}} — \emph{iḥtibās ḥarārī} — \textquotedblleft{}global warming\textquotedblright{}.

\begin{quote}
\ar{قرأت} \ar{عن} \ar{الاحتباس} \ar{الحراري} \ar{في} \ar{الجريدة}.
\end{quote}

\begin{quote}
\emph{qaraʾtu ʿan al-iḥtibās al-ḥarārī fī al-jarīda.}
\end{quote}

\begin{quote}
I read about global warming in the newspaper.
\end{quote}

\begin{cousinweb}
\textbf{\ar{احتباس} \ar{حراري}} (\emph{iḥtibās ḥarārī}) is a modern coinage built from two native Arabic words (Lane gives both roots): \emph{iḥtibās}, \textquotedblleft{}being held in\textquotedblright{}, on \emph{ḥ-b-s}, \textquotedblleft{}to confine\textquotedblright{}, and \emph{ḥarārī}, \textquotedblleft{}of heat\textquotedblright{}, on \emph{ḥ-r-r}: heat held in.
\end{cousinweb}

\begin{grammarlens}[title={What you've built}]
One more word for getting around.
\end{grammarlens}

\subsection*{Your turn}
\begin{itemize}
\raggedright
  \item \emph{Say it:} \emph{iḥtibās ḥarārī}
  \item \emph{Say it:} \emph{iḥtibās ḥarārī}, once more
  \item \emph{Say it:} \emph{qaraʾtu ʿan al-iḥtibās al-ḥarārī fī al-jarīda.} — I read about global warming in the newspaper.
  \item \emph{Recall:} say the Arabic for to raise, then the Arabic for climate, then say \emph{iḥtibās ḥarārī} again
\end{itemize}

\begin{tcolorbox}[breakable,colback=teal!4,colframe=teal!35!black,title={Before you move on}]
What does \ar{احتباس} \ar{حراري} mean? (\textquotedblleft{}Global warming\textquotedblright{}.) Say it once more.
\end{tcolorbox}
\section[\texorpdfstring{\ar{غلاف} \ar{جوي} (ghilāf jawwī)}{غلاف جوي (ghilāf jawwī)}]{\ar{غلاف} \ar{جوي} (ghilāf jawwī) — the atmosphere}

\label{lesson:AR-C463-ghilaf-jawwi}



\begin{quote}
Before the new one: say the Arabic for climate, then the Arabic for global warming.
\end{quote}

\subsection*{You'll want to know: \ar{غلاف} \ar{جوي}}
\textbf{\ar{غلاف} \ar{جوي}} — \emph{ghilāf jawwī} — \textquotedblleft{}the atmosphere\textquotedblright{}.

\begin{quote}
\ar{يحمي} \ar{الغلاف} \ar{الجوي} \ar{الأرض}.
\end{quote}

\begin{quote}
\emph{yaḥmī al-ghilāf al-jawwī al-arḍ.}
\end{quote}

\begin{quote}
The atmosphere protects the earth.
\end{quote}

\begin{cousinweb}
\textbf{\ar{غلاف} \ar{جوي}} (\emph{ghilāf jawwī}) is a modern coinage built from two native Arabic words (Lane gives both roots): \emph{ghilāf}, \textquotedblleft{}a cover, a sheath\textquotedblright{}, on \emph{gh-l-f}, and \emph{jawwī}, \textquotedblleft{}of the air\textquotedblright{}, on \emph{j-w-w}.
\end{cousinweb}

\begin{grammarlens}[title={What you've built}]
One more word for getting around.
\end{grammarlens}

\subsection*{Your turn}
\begin{itemize}
\raggedright
  \item \emph{Say it:} \emph{ghilāf jawwī}
  \item \emph{Say it:} \emph{ghilāf jawwī}, once more
  \item \emph{Say it:} \emph{yaḥmī al-ghilāf al-jawwī al-arḍ.} — The atmosphere protects the earth.
  \item \emph{Recall:} say the Arabic for climate, then the Arabic for global warming, then say \emph{ghilāf jawwī} again
\end{itemize}

\begin{tcolorbox}[breakable,colback=teal!4,colframe=teal!35!black,title={Before you move on}]
What does \ar{غلاف} \ar{جوي} mean? (\textquotedblleft{}The atmosphere\textquotedblright{}.) Say it once more.
\end{tcolorbox}
\section[\texorpdfstring{\ar{انبعاث} (inbiʿāth)}{انبعاث (inbiʿāth)}]{\ar{انبعاث} (inbiʿāth) — an emission; a release}

\label{lesson:AR-C463-inbiath}



\begin{quote}
Before the new one: say the Arabic for global warming, then the Arabic for the atmosphere.
\end{quote}

\subsection*{You'll want to know: \ar{انبعاث}}
\textbf{\ar{انبعاث}} — \emph{inbiʿāth} — \textquotedblleft{}an emission; a release\textquotedblright{}.

\begin{quote}
\ar{كان} \ar{الانبعاث} \ar{من} \ar{المصنع} \ar{كبيرا}.
\end{quote}

\begin{quote}
\emph{kāna al-inbiʿāth min al-maṣnaʿ kabīran.}
\end{quote}

\begin{quote}
The emission from the factory was large.
\end{quote}

\begin{cousinweb}
\textbf{\ar{انبعاث}} (\emph{inbiʿāth}) is built on the native Arabic root \emph{b-ʿ-th}, \textquotedblleft{}to send forth, to rouse\textquotedblright{} (Lane). It is the verbal noun of Form VII, \textquotedblleft{}being sent out\textquotedblright{}.
\end{cousinweb}

\begin{grammarlens}[title={What you've built}]
One more word for getting around.
\end{grammarlens}

\subsection*{Your turn}
\begin{itemize}
\raggedright
  \item \emph{Say it:} \emph{inbiʿāth}
  \item \emph{Say it:} \emph{inbiʿāth}, once more
  \item \emph{Say it:} \emph{kāna al-inbiʿāth min al-maṣnaʿ kabīran.} — The emission from the factory was large.
  \item \emph{Recall:} say the Arabic for global warming, then the Arabic for the atmosphere, then say \emph{inbiʿāth} again
\end{itemize}

\begin{tcolorbox}[breakable,colback=teal!4,colframe=teal!35!black,title={Before you move on}]
What does \ar{انبعاث} mean? (\textquotedblleft{}An emission; a release\textquotedblright{}.) Say it once more.
\end{tcolorbox}
\section[\texorpdfstring{\ar{كربون} (karbūn)}{كربون (karbūn)}]{\ar{كربون} (karbūn) — carbon}

\label{lesson:AR-C463-karbun}



\begin{quote}
Before the new one: say the Arabic for the atmosphere, then the Arabic for an emission.
\end{quote}

\subsection*{You'll want to know: \ar{كربون}}
\textbf{\ar{كربون}} — \emph{karbūn} — \textquotedblleft{}carbon\textquotedblright{}.

\begin{quote}
\ar{في} \ar{الهواء} \ar{كربون} \ar{كثير} \ar{فوق} \ar{المدينة}.
\end{quote}

\begin{quote}
\emph{fī al-hawāʾ karbūn kathīr fawqa al-madīna.}
\end{quote}

\begin{quote}
There is a lot of carbon in the air above the city.
\end{quote}

\begin{cousinweb}
\textbf{\ar{كربون}} (\emph{karbūn}) is usually traced to French \emph{carbone}, from Latin \emph{carbo}, \textquotedblleft{}coal\textquotedblright{}.
\end{cousinweb}

\begin{grammarlens}[title={What you've built}]
One more word for getting around.
\end{grammarlens}

\subsection*{Your turn}
\begin{itemize}
\raggedright
  \item \emph{Say it:} \emph{karbūn}
  \item \emph{Say it:} \emph{karbūn}, once more
  \item \emph{Say it:} \emph{fī al-hawāʾ karbūn kathīr fawqa al-madīna.} — There is a lot of carbon in the air above the city.
  \item \emph{Recall:} say the Arabic for the atmosphere, then the Arabic for an emission, then say \emph{karbūn} again
\end{itemize}

\begin{tcolorbox}[breakable,colback=teal!4,colframe=teal!35!black,title={Before you move on}]
What does \ar{كربون} mean? (\textquotedblleft{}Carbon\textquotedblright{}.) Say it once more.
\end{tcolorbox}
